% SPDX-FileCopyrightText: Copyright (c) 2016-2026 Objectionary.com
% SPDX-License-Identifier: MIT

\newpage
\section{Example of Dataization}
\label{app:dataization-examples}

Consider an algorithm that converts degrees from Celsius to Fahrenheit:
%
\begin{equation*}
^{\circ}F ={} ^{\circ}C \times 1.8 + 32,
\end{equation*}
%
Assuming that \(^{\circ}C\) already equals \ff{25} and is attached
  to the \ff{c} attribute instead of being provided by a user,
  the algorithm may be represented by an expression:
%
\iexec[maybe]{phino rewrite --output=latex --sweet --label=eq:celsius --margin=80 --canonize ./examples/celsius.phi > _tex/celsius-formula.tex}
\input{_tex/celsius-formula.tex}

The functions are defined as follows:
\begin{align*}
\adjustbox{max width=\linewidth,margin=0em .3em}{%
  \text{\texttt{F1}:}\;%
  \ensuremath{\infer{\text{\phiq{\phinoEvaluate{b}{e}{s}{ n_3 }{s}}}}{%
    \infer[\;\text{if}\;\text{\phiq{\delta := \delta_1 \times \delta_2}}]{\text{\phiq{\phinoNormalize{ e.number( e.bytes( data -> [[ D> \delta ]] ) ) }{ n_3 }}}}{%
      \infer{\text{\phiq{\phinoDataize{n_1}{e}{s}{\delta_1}{s}}}}{%
        \text{\phiq{\phinoNormalize{b.^}{n_1}}}%
      }%
      &
      \infer{\text{\phiq{\phinoDataize{n_2}{e}{s}{\delta_2}{s}}}}{%
        \text{\phiq{\phinoNormalize{b.|x|}{n_2}}}%
      }%
    }%
  }}%
}
\\
\adjustbox{max width=\linewidth,margin=0em .3em}{%
  \text{\texttt{F2}:}\;%
  \ensuremath{\infer{\text{\phiq{\phinoEvaluate{b}{e}{s}{ n_3 }{s}}}}{%
    \infer[\;\text{if}\;\text{\phiq{\delta := \delta_1 + \delta_2}}]{\text{\phiq{\phinoNormalize{ e.number( e.bytes( data -> [[ D> \delta ]] ) ) }{ n_3 }}}}{%
      \infer{\text{\phiq{\phinoDataize{n_1}{e}{s}{\delta_1}{s}}}}{%
        \text{\phiq{\phinoNormalize{b.^}{n_1}}}%
      }%
      &
      \infer{\text{\phiq{\phinoDataize{n_2}{e}{s}{\delta_2}{s}}}}{%
        \text{\phiq{\phinoNormalize{b.|x|}{n_2}}}%
      }%
    }%
  }}%
}
\end{align*}

The following sequence of transformations constitutes
  the dataization of \cref{eq:celsius}:

\iexec[maybe]{phino dataize --output=latex --sweet --nonumber --compress --canonize --meet-prefix=dataization --sequence --flat --quiet --hide=Q.bytes --hide=Q.number --locator=Q.@ --focus=Q.@ --meet-length=5 --meet-popularity=1 ./examples/celsius.phi > _tex/celsius-dataization.tex}
\newsavebox{\celsiusbox}
% The derivation is one tall stack of long lines, so we typeset it as an inline
%   \aligned of its natural width, rather than a page-width display, letting
%   \adjustbox rotate and scale it to the page without an overfull box.
\expandafter\let\expandafter\celsiusEq\csname equation*\endcsname
\expandafter\let\expandafter\celsiusEnd\csname endequation*\endcsname
\let\celsiusSplit\split
\let\celsiusEndsplit\endsplit
\let\celsiusGathered\gathered
\let\celsiusEndgathered\endgathered
\renewenvironment{equation*}{\(}{\)}
\renewenvironment{split}{\aligned}{\endaligned}
\renewenvironment{gathered}{\aligned}{\endaligned}
\sbox{\celsiusbox}{\input{_tex/celsius-dataization.tex}}
\expandafter\let\csname equation*\endcsname\celsiusEq
\expandafter\let\csname endequation*\endcsname\celsiusEnd
\let\split\celsiusSplit
\let\endsplit\celsiusEndsplit
\let\gathered\celsiusGathered
\let\endgathered\celsiusEndgathered
\noindent\adjustbox{angle=-90}{\adjustbox{max width=.95\textheight,max height=\linewidth}{\usebox{\celsiusbox}}}
